% !TeX root = bericht.tex
\section{Aufgabe 2}
\subsection{Vorbereitung}
Die Übertragungsfunktion eines RC-Tiefpassfilters lautet:
\begin{align}
    G(s) = \frac{1}{1 + sRC}
\end{align}
Da die Grenzfunktion asympotisch stabil ist, können wir es in den Frequenzbereich überführen:
\begin{equation}
    G(s) = \left. \frac{1}{1 + sRC} \right|_{s = j\omega} = \frac{1}{1 + j\omega RC}
\end{equation}
Die Grenzfrequenz $f_g$ eines RC-Tiefpassfilters ist definiert als die Frequenz, bei der die Amplitude des Ausgangssignals um 3dB abgefallen ist. 
Sie kann mit der folgenden Formel berechnet werden:
\begin{align}
    f_g = \frac{1}{2 \pi RC}
\end{align}
Da wir $R$ und $C$ suchen, stellen wir um. Hier wird nach $C$ umgestellt und $R$ auf einen festen Wert von $R = 10 k\Omega$ gesetzt:
\begin{align}
    C &= \frac{1}{2 \pi R f_g} \\
    &= \frac{1}{2 \pi \cdot 10 k\Omega \cdot 10 kHz} \approx 1.59 nF
\end{align}

\subsection{Grenzfrequenz in MatLab beschreiben}
\begin{sourcecode}[H]
    \caption{Matlab Code zur Beschreibung der Übertragungsfunktion}
\begin{minted}{matlab}
f_grenz = 10e3;     % 10kHz
R = 10e3;   % 10kOhm
C = 1 / (2 * pi * f_grenz * R);

num = [1];
denum = [R*C, 1];

G = tf(num, denum)
\end{minted}
\label{code:tran_fun}
\end{sourcecode}

Die Übertragungsfunktion kann in MatLab mit der Funktion \texttt{tf} beschrieben werden.

Hierbei ist auf die spezielle Syntax der Funktion \texttt{tf} zu achten.
Die \texttt{tf} Funktion erwartet zwei Arrays mit jeweils die Koeffizienten des Zählers und des Nenners der Übertragungsfunktion.
Die Koeffizienten müssen in absteigender Potenz von $s$ angegeben werden. \\
Betrachten wir als Beispiel die Übertragungsfunktion $G(s) = \frac{67s+420}{69s^2 + 42s + 1}$.
Der Zähler muss als \texttt{[67, 420]} und der Nenner als \texttt{[69, 42, 1]} angegeben werden. \\
Machen wir das für unseren RC-Tiefpass, so erhalten wir die folgenden Übertragungsfunktion:
\begin{consoleoutput}[H]
    \caption{Übertragungsfunktion des RC-Tiefpasses}
    \begin{minted}[linenos=false]{text}
    G =
 
         1
  ---------------
  1.592e-05 s + 1
    \end{minted}
\end{consoleoutput}

\subsection{Bode-Diagramm}
\begin{figure}[H]
    \centering
    \includegraphics[width=0.9\textwidth]{lab1_2_0.png}
    \caption{Bode-Diagramm des RC-Tiefpassfilters}
    \label{fig:bode}
\end{figure}

Abbildung \ref{fig:bode} zeigt das Bode-Diagramm des RC-Tiefpassfilters.
Es zeigt die Amplitude und die Phase der Übertragungsfunktion über einem Frequenzintervall von $10 Hz \le f \le 1 MHz$. 
Damit wir ein Intervall von genau $10 Hz \le f \le 1 MHz$ abbilden, müssen wir noch die Grenzen definieren mit \texttt{grenzen = \{ <untere Grenze>, <obere Grenze> \}}.
Da die Frequenz in standardmäßig in Rad/s angegeben wird, müssen wir die Option \texttt{FreqUnits} auf \texttt{Hz} setzen.

\begin{sourcecode}[H]
    \caption{Bode-Diagramm plotten}
    \begin{minted}{matlab}
grenzen = {2*pi*10, 2*pi*1e6}; % Geschweifte Klammer für Intervall
bodeOpts = bodeoptions;
bodeOpts.FreqUnits = 'Hz';
bodeOpts.Grid = 'on';   

bodeplot(G, grenzen, bodeOpts)
title("RC-Tiefpass erster Ordnung")
xlabel("Frequenz f")
    \end{minted}
\end{sourcecode}

\subsection{Nyquist-Diagramm}
Ein Nyquist-Diagramm ist, wie auch das Bode-Diagramm, eine grafische Darstellung der Übertragungsfunktion eines LTI-Systems. \\
Ein Bodediagramm zeigt die Amplitude und Phase und trägt diese jeweils über der Frequenz auf:
\begin{align}
    \text{Amplitude} &= |G(j\omega)| \nonumber \\
    \text{Phase} &= \angle{G(j\omega)}) \nonumber
\end{align}
Ein Nyquist-Diagramm zeigt die Real- und Imaginärteile der Übertragungsfunktion:
\begin{align}   
    \text{Realteil} &= \Re{G(j\omega)} \nonumber \\
    \text{Imaginärteil} &= \Im{G(j\omega)} \nonumber
\end{align}
Der Imaginärteil wird über den Realteil aufgetragen, also $\Im{G(j\omega)}$ auf der y-Achse und $\Re{G(j\omega)}$ auf der x-Achse.

In unserem Fall:
\begin{align}
    \Re{G(j\omega)} &= \frac{1}{1 + (\omega RC)^2} \nonumber \\
    \Im{G(j\omega)} &= \frac{-\omega RC}{1 + (\omega RC)^2} \nonumber
\end{align}
Man zeichnet die Kurve für $\omega$ von $0$ bis $\infty$, gekennzeichnet durch die Pfeile in der Kurve. \\

\begin{figure}[H]
    \centering
    \includegraphics[width=0.95\textwidth]{lab1_2_1.png}
    \caption{Nyquist-Diagramm des RC-Tiefpassfilters}
    \label{fig:nyquist-ellipse}
\end{figure}

\begin{sourcecode}[H]
    \caption{Nyquist-Diagramm Ellipse plotten}
    \begin{minted}{matlab}
figure;
npOpts = nyquistoptions;
nyquistplot(G, npOpts)
title("Nyquist Diagram (Ellipse)")
    \end{minted}
\end{sourcecode}

Man erkennt in Abbildung \ref{fig:nyquist-ellipse} jedoch, dass der Nyquist-Plot eine Ellipse ist,
anstatt wie üblich ein Kreis.
Dies liegt daran, dass die Achsen nicht im gleichen Maßstab skaliert sind.
Um dies zu beheben, müssen wir nur die Option \texttt{axis equal} hinzufügen.

\begin{sourcecode}[H]
    \caption{Nyquist-Diagramm als Kreis plotten}
    \begin{minted}{matlab}
figure;
npOpts = nyquistoptions;
nyquistplot(G, npOpts)
axis equal;     % <-- Achsen im gleichen Maßstab
title("Nyquist Diagram (Kreis)")
    \end{minted}
\end{sourcecode}

\begin{figure}[H]
    \centering
    \includegraphics[width=0.95\textwidth]{lab1_2_2.png}
    \caption{Nyquist-Diagramm des RC-Tiefpassfilters (Kreis)}
    \label{fig:nyquist-kreis}
\end{figure}

Uns ist jedoch bekannt, dass die Übertragungsfunktion eines RC-Tiefpassfilters ein halbkreis 
im vierten Quadranten ist, jedoch wird in Abbildung \ref{fig:nyquist-kreis} ein voller Kreis dargestellt.
Dies liegt daran, dass die Nyquist-Plot Funktion von MatLab die Frequenz von $-\infty$ bis $\infty$ abbildet, anstatt von $0$ bis $\infty$.
Um dies zu beheben, müssen wir die Option \texttt{npOpts.ShowFullContour = 'off';} hinzufügen.
\begin{sourcecode}[H]
    \caption{Nyquist-Diagramm als Halbkreis plotten}
    \begin{minted}{matlab}
f = figure;
npOpts.ShowFullContour = 'off';     % <-- exclude negative frequencies
nyquistplot(G, npOpts)
axis equal;     % <-- Damit es als Hlabkreis dargestellt wird und nicht als halbellipse
title("Nyquist Diagram (Halbkreis)")
    \end{minted}
\end{sourcecode}

\begin{figure}[H]
    \centering
    \includegraphics[width=0.9\textwidth]{lab1_2_3.png}
    \caption{Nyquist-Diagramm des RC-Tiefpassfilters (Halbkreis)}
    \label{fig:nyquist-halbkreis}
\end{figure}
